\addtocontents{toc}{\protect\Needspace{3\baselineskip}}
\originalchapter{弦、紧致维度与作用量}\label{ch:action}
\originalsection{从粒子轨迹到世界面}
点粒子的历史是一条世界线, 弦的每个空间点随时间运动后形成二维世界面. 弦的尺度不能通过把点粒子方程简单增加一个坐标得到; 需要说明弦上的参数怎样标记同一个物理对象, 哪些变形是实际运动, 哪些仅是重新命名.

目标时空坐标 $X^\mu(\tau,\sigma)$ 是世界面上的标量场, $\mu=0,\ldots,D-1$ 是目标时空指标, $a,b=0,1$ 是世界面指标. 点粒子的相对论作用量 $S=-m\int\sqrt{-\eta_{\mu\nu}\dot X^\mu\dot X^\nu}\dd\tau$ 只依固有长度, 在 $\tau$ 的单调重参数化下不变. 弦则以世界面面积给
\begin{equation}\label{eq:NG}
S_{\rm NG}=-T_s\int\dd^2\sigma\sqrt{-\det h_{ab}},\qquad
h_{ab}=\eta_{\mu\nu}\partial_aX^\mu\partial_bX^\nu,
\quad T_s=\frac1{2\pi\alpha'}.
\end{equation}
诱导度规 $h$ 不是目标时空度规本身, 它把目标空间内积拉回到世界面. $\alpha'$ 在自然单位下为长度平方, 控制弦长与质量尺度. 面積的不变性使二维重参数化成为规范自由度.

非相对论横向小振动作用量是 $S=\frac12\int\dd t\dd x[\mu_0\dot y^2-T_0(y')^2]$. 变分得 $\mu_0\ddot y-T_0y''=0$, 波速 $\sqrt{T_0/\mu_0}$. Nambu--Goto 作用量在静态规范 $X^0=ct,X^1=x$ 下展开, 给同一二次形式且 $\mu_0=T_0/c^2$, 说明相对论弦的横波速度是 $c$. 仅当横向速度和斜率都小时才允许此展开.

\originalsection{Polyakov 形式与约束}
面积平方根不便量子化. 引入独立世界面度规 $\gamma_{ab}$, 取
\begin{equation}\label{eq:Polyakov}
S_P=-\frac1{4\pi\alpha'}\int\dd^2\sigma\sqrt{-\gamma}\gamma^{ab}
\partial_aX^\mu\partial_bX_\mu.
\end{equation}
对 $\gamma^{ab}$ 变分给 $T_{ab}=\partial_aX\cdot\partial_bX-\gamma_{ab}\gamma^{cd}\partial_cX\cdot\partial_dX/2=0$. 在非退化世界面上, 此式使 $\gamma$ 与 $h$ 仅差共形因子, 代回恢复式\eqref{eq:NG}. 二维作用量在 $\gamma_{ab}\to\ee^{2\omega}\gamma_{ab}$ 下不变, 因 $\sqrt{-\gamma}$ 与 $\gamma^{ab}$ 的缩放正好相消.

局部可选共形规范 $\gamma_{ab}=\ee^{2\omega}\eta_{ab}$, 自由平直目标空间的场方程变为 $(\partial_\tau^2-\partial_\sigma^2)X^\mu=0$. 但度规方程仍须施加, 即
\begin{equation}\label{eq:constraints}
\dot X^2+X'^2=0,\qquad\dot X\cdot X'=0,
\quad(\dot X\pm X')^2=0.
\end{equation}
「场方程线性」不意味着任意线性解都是物理弦. 忽略约束会把规范方向误当作独立振子, 得到错误状态数.

\originalsection{紧致化与商空间}
若空间方向 $y\sim y+2\pi R$, 单值波函数 $\ee^{\ii p_y y}$ 要求 $p_y=n/R$, 低维观察者于是看到 Kaluza--Klein 质量 $m_n^2=m_0^2+n^2/R^2$. 闭弦还可绕圆 $w$ 次, 最小伸长能对应质量尺度 $|w|R/\alpha'$, 故谱可同时含动量和绕数. 把 $R\leftrightarrow\alpha'/R,n\leftrightarrow w$ 交换可提示 T 对偶, 但完整相互作用与背景下的对偶证明不在此入门推导中.

平面商 $z\sim\ee^{2\pi\ii/N}z$ 的基本域是角度 $2\pi/N$ 的扇形, 两边须识别. 圆锥尖点有全局几何信息, 平面上普通直线在商空间中可能需要跨域拼接. 基本域必须每轨道仅取一次代表, 不能因为生成元旋转 $2\pi M/N$ 就把整片较大的扇形当成基本域.

\originalsection{面积与小振动的逐项展开}
静态规范取 $X^0=ct,X^1=x,\boldsymbol X_\perp=\boldsymbol y(t,x)$. 诱导度规分量为 $h_{tt}=-c^2+\dot{\boldsymbol y}^2$, $h_{xx}=1+\boldsymbol y'^2$, $h_{tx}=\dot{\boldsymbol y}\cdot\boldsymbol y'$. 恢复物理张力 $T_0$ 后,
\[
S=-T_0\int\dd t\dd x\sqrt{1+\boldsymbol y'^2-
\frac{\dot{\boldsymbol y}^2}{c^2}
-\frac{\dot{\boldsymbol y}^2\boldsymbol y'^2-(\dot{\boldsymbol y}\cdot\boldsymbol y')^2}{c^2}}.
\]
最后一项是四阶小量. 对小斜率与小速度同时展开平方根, 得 $S=-T_0\int\dd t\dd x+\frac12\int\dd t\dd x[(T_0/c^2)\dot y^2-T_0y'^2]+O(y^4)$. 因而基态静能是 $T_0$ 乘长度, 动能系数的线密度为 $T_0/c^2$. 若仅速度小而弦弯曲很大, 不能把所有 $y'^2$ 都作为小量.

Polyakov 作用量的度规变化也可明确复算. $\delta\sqrt{-\gamma}=-\frac12\sqrt{-\gamma}\gamma_{ab}\delta\gamma^{ab}$, 于是 $\delta S_P$ 的系数为 $h_{ab}-\gamma_{ab}\gamma^{cd}h_{cd}/2$. 零应力意味着 $h_{ab}=\lambda\gamma_{ab}$, $\lambda=\gamma^{cd}h_{cd}/2$. 在物理非退化支, $\sqrt{-h}=\lambda\sqrt{-\gamma}$, 代回 $S_P=-T_s\int\sqrt{-h}$. 用这些对象类型可以避免把 $\gamma$ 固定以后忘记保留它原来的方程.
